\chapter{2017年北京卷（理）}

\section{单选题}

\begin{problem}
若集合 \(A=\{x\mid -2<x<1\}\)，\(B=\{x\mid x<-1\text{ 或 }x>3\}\)，则 \(A\cap B=\)（\quad）

\choices
    {\(\{x\mid -2<x<-1\}\)}
    {\(\{x\mid -2<x<3\}\)}
    {\(\{x\mid -1<x<1\}\)}
    {\(\{x\mid 1<x<3\}\)}

\end{problem}

\begin{answer}
A
\end{answer}

\begin{solution}
由题意，\(A\cap B\) 中的元素同时满足 \(-2<x<1\) 和 \(x<-1\text{ 或 }x>3\)，因此只能满足 \(-2<x<-1\). 故选 \(A\).
\end{solution}
\begin{problem}
若复数 \((1-\i)(a+\i)\) 在复平面内对应的点在第二象限，则实数 \(a\) 的取值范围是（\quad）

\choices
    {\((-\infty,1)\)}
    {\((-\infty,-1)\)}
    {\((1,+\infty)\)}
    {\((-1,+\infty)\)}

\end{problem}

\begin{answer}
B
\end{answer}

\begin{solution}
将复数化简得
\((1-\i)(a+\i)=a+1+(1-a)\i\)
要使其对应的点在第二象限，需实部小于 \(0\)、虚部大于 \(0\)，即
\(a+1<0\)，\(1-a>0\)
由此得 \(a<-1\). 故选 \(B\).
\end{solution}
\begin{problem}
执行如图所示的程序框图，输出的 \(s\) 值为（\quad）

\begin{center}
\bitmapinclude[max width=0.42\linewidth,max totalheight=4.8cm]{img/2017/beijing_science/q3_fig1.png}
\end{center}

\choices
    {\(2\)}
    {\(\dfrac32\)}
    {\(\dfrac53\)}
    {\(\dfrac85\)}

\end{problem}

\begin{answer}
C
\end{answer}

\begin{solution}
按程序框图依次执行，最终输出的 \(s\) 为 \(\dfrac53\). 故选 \(C\).
\end{solution}
\begin{problem}
若 \(x,y\) 满足约束条件 \(x\le 3\)，\(x+y\ge 2\)，\(y\le x\)，则 \(x+2y\) 的最大值为（\quad）

\choices
    {\(1\)}
    {\(3\)}
    {\(5\)}
    {\(9\)}

\end{problem}

\begin{answer}
D
\end{answer}

\begin{solution}
由约束条件作出可行域，其三个顶点为 \((3,-1)\)、\((3,3)\)、\((1,1)\). 分别计算 \(x+2y\) 得
\(3+2(-1)=1\)，\(3+2\cdot 3=9\)，\(1+2\cdot 1=3\)
故最大值为 \(9\)，选 \(D\).
\end{solution}
\begin{problem}
已知函数 \(f(x)=3^x-\left(\dfrac13\right)^x\)，则 \(f(x)\)（\quad）

\choices
    {是奇函数，且在 \(\R\) 上是增函数}
    {是偶函数，且在 \(\R\) 上是增函数}
    {是奇函数，且在 \(\R\) 上是减函数}
    {是偶函数，且在 \(\R\) 上是减函数}

\end{problem}

\begin{answer}
A
\end{answer}

\begin{solution}
因为
\(f(-x)=3^{-x}-\left(\dfrac13\right)^{-x}=\left(\dfrac13\right)^x-3^x=-f(x)\)
所以 \(f(x)\) 是奇函数. 又
\(f'(x)=3^x\ln 3+\left(\dfrac13\right)^x\ln 3>0\)
所以 \(f(x)\) 在 \(\R\) 上是增函数. 故选 \(A\).
\end{solution}
\begin{problem}
设 \(\bs{m},\bs{n}\) 为非零向量，则“存在负数 \(\lambda\)，使得 \(\bs{m}=\lambda\bs{n}\)”是“\(\bs{m}\cdot \bs{n}<0\)”的（\quad）

\choices
    {充分而不必要条件}
    {必要而不充分条件}
    {充分必要条件}
    {既不充分也不必要条件}

\end{problem}

\begin{answer}
A
\end{answer}

\begin{solution}
若存在负数 \(\lambda\)，使得 \(\bs{m}=\lambda\bs{n}\)，则
\(\bs{m}\cdot \bs{n}=\lambda\lVert \bs{n}\rVert^2<0\)
所以这是充分条件. 反之，\(\bs{m}\cdot \bs{n}<0\) 只能说明两向量夹角为钝角，不一定共线，所以不是必要条件. 故选 \(A\).
\end{solution}
\begin{problem}
某四棱锥的三视图如图所示，则该四棱锥的最长棱的长度为（\quad）

\begin{center}
\bitmapinclude[max width=0.38\linewidth,max totalheight=4.8cm]{img/2017/beijing_science/q7_fig1.png}
\end{center}

\choices
    {\(3\sqrt2\)}
    {\(2\sqrt3\)}
    {\(2\sqrt2\)}
    {\(2\)}

\end{problem}

\begin{answer}
B
\end{answer}

\begin{solution}
由三视图还原该四棱锥，最长棱为相应空间对角线，长度为 \(2\sqrt3\). 故选 \(B\).
\end{solution}
\begin{problem}
根据有关资料，围棋状态空间复杂度的上限 \(M\) 约为 \(3^{361}\)，而可观测宇宙中普通物质的原子总数 \(N\) 约为 \(10^{80}\). 则下列各数中与 \(\dfrac{M}{N}\) 最接近的是（\quad）

（参考数据：\(\lg 3\approx 0.48\)）

\choices
    {\(10^{33}\)}
    {\(10^{53}\)}
    {\(10^{73}\)}
    {\(10^{93}\)}

\end{problem}

\begin{answer}
D
\end{answer}

\begin{solution}
有
\(\lg\frac{M}{N}=361\lg 3-80\approx 361\times 0.48-80=93.28\)
所以 \(\dfrac{M}{N}\) 约为 \(10^{93}\). 故选 \(D\).
\end{solution}
\section{填空题}

\begin{problem}
若双曲线 \(\dfrac{x^2}{1}-\dfrac{y^2}{m}=1\) 的离心率为 \(\sqrt3\)，则实数 \(m=\fillinblank{}\).

\end{problem}

\begin{answer}
2
\end{answer}

\begin{solution}
双曲线的离心率
\(e=\sqrt{1+\frac{m}{1}}=\sqrt{1+m}\)
由 \(e=\sqrt3\) 得 \(1+m=3\)，所以 \(m=2\).
\end{solution}
\begin{problem}
若等差数列 \(\{a_n\}\) 和等比数列 \(\{b_n\}\) 满足 \(a_1=b_1=-1\)，\(a_4=b_4=8\)，则 \(\dfrac{a_2}{b_2}=\fillinblank{}\).

\end{problem}

\begin{answer}
1
\end{answer}

\begin{solution}
设等差数列 \(\{a_n\}\) 的公差为 \(d\)，则
\(a_4=a_1+3d=-1+3d=8\)
所以 \(d=3\)，从而 \(a_2=2\).
设等比数列 \(\{b_n\}\) 的公比为 \(q\)，则
\(b_4=b_1q^3=-q^3=8\)
所以 \(q=-2\)，从而 \(b_2=2\).
故 \(\dfrac{a_2}{b_2}=1\).
\end{solution}
\begin{problem}
在极坐标系中，点 \(A\) 在圆 \(\rho^2-2\rho\cos\theta-4\rho\sin\theta+4=0\) 上，点 \(P\) 的坐标为 \((1,0)\)，则 \(|AP|\) 的最小值为 \fillinblank{}.

\end{problem}

\begin{answer}
1
\end{answer}

\begin{solution}
将极坐标方程化为直角坐标方程得
\(x^2+y^2-2x-4y+4=0\)
即
\((x-1)^2+(y-2)^2=1\)
所以圆心为 \((1,2)\)，半径为 \(1\). 点 \(P(1,0)\) 到圆心的距离为 \(2\)，故 \(|AP|\) 的最小值为 \(2-1=1\).
\end{solution}
\begin{problem}
在平面直角坐标系 \(xOy\) 中，角 \(\alpha\) 与角 \(\beta\) 均以 \(Ox\) 为始边，它们的终边关于 \(y\) 轴对称. 若 \(\sin\alpha=\dfrac13\)，则 \(\cos(\alpha-\beta)=\fillinblank{}\).

\end{problem}

\begin{answer}
\(-\dfrac79\)
\end{answer}

\begin{solution}
由题意知 \(\beta=\pi-\alpha\)，所以
\(\cos(\alpha-\beta)=\cos(2\alpha-\pi)=-\cos 2\alpha\)
又因为 \(\sin\alpha=\dfrac13\)，所以
\(\cos 2\alpha=1-2\sin^2\alpha=1-\frac29=\frac79\)
故 \(\cos(\alpha-\beta)=-\dfrac79\).
\end{solution}
\begin{problem}
能够说明“设 \(a,b,c\) 是任意实数，若 \(a>b>c\)，则 \(a+b>c\)”是假命题的一组整数 \(a,b,c\) 的值依次为 \fillinblank{}.

\end{problem}

\begin{answer}
\(-1,-2,-3\)
\end{answer}

\begin{solution}
取 \(a=-1\)，\(b=-2\)，\(c=-3\)，则 \(a>b>c\) 成立，但
\(a+b=-3=c\)
所以 \(a+b>c\) 不成立. 故这是一组可以说明原命题是假命题的整数.
\end{solution}
\begin{problem}
三名工人加工同一种零件，他们在一天中的工作情况如图所示，其中点 \(A_i\) 的横、纵坐标分别为第 \(i\) 名工人上午的工作时间和加工的零件数，点 \(B_i\) 的横、纵坐标分别为第 \(i\) 名工人下午的工作时间和加工的零件数，\(i=1,2,3\).

\begin{center}
\bitmapinclude[max width=0.44\linewidth,max totalheight=4.8cm]{img/2017/beijing_science/q14_fig1.png}
\end{center}

\begin{circlelist}
\item 记 \(Q_i\) 为第 \(i\) 名工人在这一天中加工的零件总数，则 \(Q_1,Q_2,Q_3\) 中最大的是 \fillinblank{}.
\item 记 \(p_i\) 为第 \(i\) 名工人在这一天中平均每小时加工的零件数，则 \(p_1,p_2,p_3\) 中最大的是 \fillinblank{}.
\end{circlelist}

\end{problem}

\begin{answer}
\(Q_1,\ p_2\)
\end{answer}

\begin{solution}
由图像可比较三名工人上午与下午加工零件数之和，得 \(Q_1\) 最大. 再比较“总加工数除以总工作时长”，故 \(p_2\) 最大.
\end{solution}
\section{解答题}

\begin{problem}
在 \(\triangle ABC\) 中，\(\angle A=60^\circ\)，\(c=\dfrac37a\).

\begin{enumerate}
\item 求 \(\sin C\) 的值；
\item 若 \(a=7\)，求 \(\triangle ABC\) 的面积.
\end{enumerate}

\end{problem}

\begin{solution}
（1）由正弦定理得
\(\frac{c}{\sin C}=\frac{a}{\sin A}\)
因为 \(\angle A=60^\circ\)，\(c=\dfrac37a\)，所以
\(\sin C=\frac{c\sin A}{a}=\frac37\cdot\frac{\sqrt3}{2}=\frac{3\sqrt3}{14}\)

（2）当 \(a=7\) 时，\(c=3\). 由余弦定理得
\(7^2=b^2+3^2-2\cdot b\cdot 3\cdot \cos60^\circ\)
即
\(49=b^2+9-3b\)
解得 \(b=8\)（舍去 \(b=-5\)）. 因此
\(S_{\triangle ABC}=\frac12bc\sin A=\frac12\cdot 8\cdot 3\cdot \frac{\sqrt3}{2}=6\sqrt3\)
\end{solution}
\begin{problem}
如图，在四棱锥 \(P-ABCD\) 中，底面 \(ABCD\) 为正方形，平面 \(PAD\perp\) 平面 \(ABCD\)，点 \(M\) 在线段 \(PB\) 上，\(PD//\) 平面 \(MAC\)，\(PA=PD=\sqrt6\)，\(AB=4\).

\begin{center}
\bitmapinclude[max width=0.52\linewidth,max totalheight=4.8cm]{img/2017/beijing_science/q16_fig1.png}
\end{center}

\begin{enumerate}
\item 求证：\(M\) 为 \(PB\) 的中点；
\item 求二面角 \(B-PD-A\) 的大小；
\item 求直线 \(MC\) 与平面 \(BDP\) 所成角的正弦值.
\end{enumerate}

\end{problem}

\begin{solution}
（1）设 \(AC\) 与 \(BD\) 交于点 \(E\)，连接 \(ME\). 因为 \(PD//\) 平面 \(MAC\)，且平面 \(MAC\cap\) 平面 \(PBD=ME\)，所以 \(PD//ME\). 又因为 \(ABCD\) 是正方形，所以 \(E\) 为 \(BD\) 的中点. 因此 \(M\) 为 \(PB\) 的中点.

（2）取 \(AD\) 的中点 \(O\)，连接 \(OP\)、\(OE\). 因为 \(PA=PD\)，所以 \(OP\perp AD\). 又因为平面 \(PAD\perp\) 平面 \(ABCD\)，且 \(OP\subset\) 平面 \(PAD\)，所以 \(OP\perp\) 平面 \(ABCD\). 因为 \(OE\subset\) 平面 \(ABCD\)，所以 \(OP\perp OE\). 又因为 \(ABCD\) 是正方形，所以 \(OE\perp AD\).

建立空间直角坐标系 \(O-xyz\)，则 \(P(0,0,\sqrt2)\)，\(D(2,0,0)\)，\(B(-2,4,0)\).
于是 \(\overrightarrow{BD}=(4,-4,0)\)，\(\overrightarrow{PD}=(2,0,-\sqrt2)\).
设平面 \(BDP\) 的法向量为 \(\bs{n}=(x,y,z)\)，则
\examdisplaycases{}{
\bs{n}\cdot \overrightarrow{BD}=0,\\
\bs{n}\cdot \overrightarrow{PD}=0,
}
即
\examdisplaycases{}{
4x-4y=0,\\
2x-\sqrt2\,z=0.
}
令 \(x=1\)，得 \(y=1\)，\(z=\sqrt2\)，故 \(\bs{n}=(1,1,\sqrt2)\).
平面 \(PAD\) 的法向量可取为 \(\bs{p}=(0,1,0)\)，于是 \(\cos\langle \bs{n},\bs{p}\rangle=\frac{\bs{n}\cdot \bs{p}}{|\bs{n}||\bs{p}|}=\frac12\).
由题知二面角 \(B-PD-A\) 为锐角，所以它的大小为 \(\dfrac{\pi}{3}\).

（3）由题意知 \(M\left(-1,2,\frac{\sqrt2}{2}\right)\)，\(C(2,4,0)\)，故 \(\overrightarrow{MC}=\left(3,2,-\frac{\sqrt2}{2}\right)\).
设直线 \(MC\) 与平面 \(BDP\) 所成角为 \(\alpha\)，则 \(\sin\alpha=\left|\cos\langle \bs{n},\overrightarrow{MC}\rangle\right|=\frac{|\bs{n}\cdot \overrightarrow{MC}|}{|\bs{n}|\,|\overrightarrow{MC}|}\).
计算得 \(\bs{n}\cdot \overrightarrow{MC}=3+2-1=4\)，\(|\bs{n}|=2\sqrt2\)，\(|\overrightarrow{MC}|=\frac{3\sqrt6}{2}\)，所以 \(\sin\alpha=\frac{4}{2\sqrt2\cdot \frac{3\sqrt6}{2}}=\frac{2\sqrt6}{9}\).
故直线 \(MC\) 与平面 \(BDP\) 所成角的正弦值为 \(\dfrac{2\sqrt6}{9}\).
\end{solution}
\begin{problem}
为了研究一种新药的疗效，选 \(100\) 名患者随机分成两组，每组各 \(50\) 名，一组服药，另一组不服药. 一段时间后，记录了两组患者的生理指标 \(x\) 和 \(y\) 的数据，并制成下图，其中“\(*\)”表示服药者，“\(+\)”表示未服药者.

\bitmapfigure[max width=0.7\linewidth,max totalheight=4.8cm]{img/2017/beijing_science/q17_fig1.png}

\begin{enumerate}
\item 从服药的 \(50\) 名患者中随机选出一人，求此人指标 \(y\) 的值小于 \(60\) 的概率；
\item 从图中 \(A,B,C,D\) 四人中随机选出两人，记 \(\xi\) 为选出的两人中指标 \(x\) 的值大于 \(1.7\) 的人数，求 \(\xi\) 的分布列和数学期望 \(E(\xi)\)；
\item 试判断这 \(100\) 名患者中服药者指标 \(y\) 数据的方差与未服药者指标 \(y\) 数据的方差的大小. （只需写出结论）
\end{enumerate}

\end{problem}

\begin{solution}
（1）由图知，在服药的 \(50\) 名患者中，指标 \(y\) 的值小于 \(60\) 的有 \(15\) 人，所以所求概率为 \(\frac{15}{50}=0.3\).

（2）由图知，\(A,B,C,D\) 四人中，指标 \(x\) 的值大于 \(1.7\) 的有 \(2\) 人：\(A\) 和 \(C\). 因此 \(\xi\) 的所有可能取值为 \(0,1,2\).

\(P(\xi=0)=\frac{\mathrm C_2^2}{\mathrm C_4^2}=\frac16\)，\(P(\xi=1)=\frac{\mathrm C_2^1\mathrm C_2^1}{\mathrm C_4^2}=\frac23\)，\(P(\xi=2)=\frac{\mathrm C_2^2}{\mathrm C_4^2}=\frac16\)

\examdisplayarray{}{c|ccc}{
\xi & 0 & 1 & 2\\
\hline
P & \dfrac16 & \dfrac23 & \dfrac16
}{}

所以 \(E(\xi)=0\cdot \frac16+1\cdot \frac23+2\cdot \frac16=1\)

（3）在这 \(100\) 名患者中，服药者指标 \(y\) 数据的方差大于未服药者指标 \(y\) 数据的方差.
\end{solution}
\begin{problem}
已知抛物线 \(C:y^2=2px\) 过点 \(P(1,1)\). 过点 \(\left(0,\dfrac12\right)\) 作直线 \(l\) 与抛物线 \(C\) 交于不同的两点 \(M,N\)，过点 \(M\) 作 \(x\) 轴的垂线分别与直线 \(OP\)、\(ON\) 交于点 \(A,B\)，其中 \(O\) 为原点.

\begin{enumerate}
\item 求抛物线 \(C\) 的方程，并求其焦点坐标和准线方程；
\item 求证：\(A\) 为线段 \(BM\) 的中点.
\end{enumerate}

\end{problem}

\begin{solution}
（1）由抛物线 \(C:y^2=2px\) 过点 \(P(1,1)\)，得 \(1=2p\)，所以 \(p=\dfrac12\)，从而抛物线 \(C\) 的方程为 \(y^2=x\)，其焦点坐标为 \(\left(\dfrac14,0\right)\)，准线方程为 \(x=-\dfrac14\).

（2）设直线 \(l\) 的方程为 \(y=kx+\frac12\)（\(k\ne 0\)），与抛物线 \(C\) 的交点为 \(M(x_1,y_1)\)、\(N(x_2,y_2)\).
由
\examdisplaycases{}{
y=kx+\dfrac12,\\
y^2=x,
}
得 \(4k^2x^2+(4k-4)x+1=0\)，所以 \(x_1+x_2=\frac{1-k}{k^2}\)，\(x_1x_2=\frac{1}{4k^2}\).
因为点 \(P\) 的坐标为 \((1,1)\)，所以直线 \(OP\) 的方程为 \(y=x\)，点 \(A\) 的坐标为 \((x_1,x_1)\).
又因为直线 \(ON\) 的方程为 \(y=\dfrac{y_2}{x_2}x\)，所以点 \(B\) 的坐标为 \(\left(x_1,\dfrac{y_2y_1}{x_2}\right)\).

由 \(y_1=kx_1+\dfrac12\)、\(y_2=kx_2+\dfrac12\)，得
\(y_1+\frac{y_2y_1}{x_2}-2x_1=\frac{y_1y_2+y_2y_1-2x_1x_2}{x_2}=\frac{\left(kx_1+\dfrac12\right)x_2+\left(kx_2+\dfrac12\right)x_1-2x_1x_2}{x_2}=\frac{(2k-2)x_1x_2+\dfrac12(x_2+x_1)}{x_2}=\frac{(2k-2)\cdot \dfrac{1}{4k^2}+\dfrac{1-k}{2k^2}}{x_2}=0\).
所以 \(y_1+\frac{y_2y_1}{x_2}-2x_1=0\). 即 \(y_1+\frac{y_2y_1}{x_2}=2x_1\).
故 \(A\) 为线段 \(BM\) 的中点.
\end{solution}
\begin{problem}
已知函数 \(f(x)=\e^x\cos x-x\).

\begin{enumerate}
\item 求曲线 \(y=f(x)\) 在点 \((0,f(0))\) 处的切线方程；
\item 求函数 \(f(x)\) 在区间 \(\left[0,\dfrac{\pi}{2}\right]\) 上的最大值和最小值.
\end{enumerate}

\end{problem}

\begin{solution}
（1）因为 \(f(x)=\e^x\cos x-x\)，所以 \(f'(x)=\e^x(\cos x-\sin x)-1\)，\(f'(0)=0\). 又因为 \(f(0)=1\)，所以曲线 \(y=f(x)\) 在点 \((0,f(0))\) 处的切线方程为 \(y=1\).

（2）设 \(h(x)=f'(x)=\e^x(\cos x-\sin x)-1\)，则 \(h'(x)=-2\e^x\sin x\).
当 \(x\in\left(0,\dfrac{\pi}{2}\right)\) 时，\(h'(x)<0\)，所以 \(h(x)\) 在区间 \(\left[0,\dfrac{\pi}{2}\right]\) 上单调递减.
因此对任意 \(x\in\left(0,\dfrac{\pi}{2}\right]\) 有 \(h(x)<h(0)=0\)，即 \(f'(x)<0\).
所以函数 \(f(x)\) 在区间 \(\left[0,\dfrac{\pi}{2}\right]\) 上单调递减.
因此最大值为 \(f(0)=1\)，最小值为 \(f\left(\dfrac{\pi}{2}\right)=\e^{\pi/2}\cdot 0-\frac{\pi}{2}=-\frac{\pi}{2}\)
\end{solution}
\begin{problem}
设 \(\{a_n\}\) 和 \(\{b_n\}\) 是两个等差数列，记 \(c_n=\max\{b_1-na_1,b_2-na_2,\dots,b_n-na_n\}\)（\(n=1,2,3,\dots\)），其中 \(\max\{x_1,x_2,\dots,x_s\}\) 表示 \(x_1,x_2,\dots,x_s\) 这 \(s\) 个数中最大的数.

\begin{enumerate}
\item 若 \(a_n=n\)，\(b_n=2n-1\)，求 \(c_1,c_2,c_3\) 的值，并证明 \(\{c_n\}\) 是等差数列；
\item 证明：或者对任意正数 \(M\)，存在正整数 \(m\)，当 \(n\ge m\) 时，\(\dfrac{c_n}{n}>M\)；或者存在正整数 \(m\)，使得 \(c_m,c_{m+1},c_{m+2},\dots\) 是等差数列.
\end{enumerate}

\end{problem}

\begin{solution}
（1）因为 \(a_n=n\)，\(b_n=2n-1\)，所以
\(c_1=b_1-a_1=1-1=0\)
\(c_2=\max\{b_1-2a_1,b_2-2a_2\}=\max\{1-2,3-4\}=-1\)
\(c_3=\max\{b_1-3a_1,b_2-3a_2,b_3-3a_3\}=\max\{1-3,3-6,5-9\}=-2\)
当 \(n\ge 3\) 时，
\((b_{k+1}-na_{k+1})-(b_k-na_k)=(b_{k+1}-b_k)-n(a_{k+1}-a_k)=2-n<0\)
所以 \(b_k-na_k\) 关于 \(k\in\N^*\) 单调递减.
因此
\(c_n=\max\{b_1-na_1,b_2-na_2,\dots,b_n-na_n\}=b_1-na_1=1-n\)
所以对任意 \(n\ge1\)，\(c_n=1-n\)，于是
\(c_{n+1}-c_n=-1\)
故 \(\{c_n\}\) 是等差数列.

（2）设数列 \(\{a_n\}\) 和 \(\{b_n\}\) 的公差分别为 \(d_1,d_2\)，则
\(b_k-na_k=b_1+(k-1)d_2-[a_1+(k-1)d_1]n=b_1-a_1n+(d_2-nd_1)(k-1)\)
所以
\examdisplaycases{c_n=}{
b_1-a_1n+(n-1)(d_2-nd_1), & d_2>nd_1,\\
b_1-a_1n, & d_2\le nd_1.
}

\circled{1}当 \(d_1>0\) 时，取正整数 \(m>\dfrac{d_2}{d_1}\)，则当 \(n\ge m\) 时，\(nd_1>d_2\)，因此 \(c_n=b_1-a_1n\).
此时，\(c_m,c_{m+1},c_{m+2},\dots\) 是等差数列.

\circled{2}当 \(d_1=0\) 时，对任意 \(n\ge1\)，\(c_n=b_1-a_1n+(n-1)\max\{d_2,0\}=b_1-a_1+(n-1)\bigl(\max\{d_2,0\}-a_1\bigr)\).
此时，\(c_1,c_2,c_3,\dots,c_n,\dots\) 是等差数列.

\circled{3}当 \(d_1<0\) 时，当 \(n>\dfrac{d_2}{d_1}\) 时，有 \(nd_1<d_2\)，所以 \(\frac{c_n}{n}=\frac{b_1-a_1n+(n-1)(d_2-nd_1)}{n}=n(-d_1)+d_1-a_1+d_2+\frac{b_1-d_2}{n}\).
从而 \(\frac{c_n}{n}\ge n(-d_1)+d_1-a_1+d_2-|b_1-d_2|\).
对任意正数 \(M\)，取正整数 \(m>\max\left\{\frac{M+|b_1-d_2|+a_1-d_1-d_2}{-d_1},\frac{d_2}{d_1}\right\}\)
故当 \(n\ge m\) 时，\(\dfrac{c_n}{n}>M\).
\end{solution}
